\documentclass{article}
\usepackage[margin=1in]{geometry}
\usepackage{amsmath,amssymb,amsthm}
\usepackage{booktabs}
\usepackage{hyperref}
\usepackage{graphicx}

\newtheorem{theorem}{Theorem}
\newtheorem{lemma}[theorem]{Lemma}
\newtheorem{corollary}[theorem]{Corollary}
\newtheorem{proposition}[theorem]{Proposition}

\title{No-Regret is Not No-Harm:\\ Bandits with Consumable Action Sets}
\author{Navikkumar Modi}
\date{}

\begin{document}
\maketitle

\begin{abstract}
Sublinear regret is the standard justification for deploying sequential
decision-making systems. We show that this justification is vacuous when actions
consume the action set. We study bandits in which pulling an arm may destroy it,
and in which each destroyed arm imposes a permanent negative externality on every
subsequent reward. We prove that a policy achieving \emph{exactly zero} regret
against the best-available-arm benchmark can incur value loss $m\delta E$, unbounded
in both the pool size $m$ and the externality magnitude $\delta E$, and can drive
total value negative while the optimal policy remains positive. The mechanism is
that regret is measured against an optimum the learner has itself degraded. We
further establish that the parameter governing optimal play, the expected marginal
externality $\kappa_a = p_a e_a$, is (i) exactly the quantity that determines
whether greedy is optimal, and (ii) subject to an estimation error floor
$2\sigma/\sqrt{\min(T,\sum_j 1/p_j)}$ that is \emph{independent of the horizon},
because each arm supplies exactly one before/after transition and the sample budget
is capped by consumption rather than time. The quantity that matters most becomes
least estimable exactly as it becomes more important. We give a closed-form
correction, characterise an exactly-solvable special case, and release an
open-source testbed.
\end{abstract}

\section{Introduction}

Consider an orchestrator routing tasks across API-backed tools. Some tools are
robust; others trip a rate limit, exhaust a quota, or get a key revoked under heavy
use, disappearing for the remainder of the session. Tools are not independent:
those sharing upstream infrastructure---a vendor, an IP pool, a quota
tier---degrade together when one is burned.

Table~\ref{tab:trace} shows two routers on identical instances. The first has a
flawless record: at every step it pulls the highest-value available tool, so its
regret against the best-available benchmark is identically zero. It ends on a
fallback tool worth $0.470$ per call with two tools remaining. The second reports
regret of $0.25$ per step---and is sitting at $0.782$ per call, $66\%$ higher, with
four tools remaining.

\begin{table}[t]
\centering
\caption{Two routers, identical seed and tools. The policy that looks worse by the
standard metric achieves the better outcome.}
\label{tab:trace}
\begin{tabular}{llrrr}
\toprule
router & step & realised & regret & tools left \\
\midrule
greedy & 0 & 1.150 & 0.0000 & 6 \\
greedy & 2 & 0.844 & 0.0000 & 4 \\
greedy & 7--13 & 0.470 & \textbf{0.0000} & \textbf{2} \\
\midrule
ECI & 0 & 1.150 & 0.0000 & 6 \\
ECI & 1--4 & 0.988 & 0.0500 & 5 \\
ECI & 5--13 & \textbf{0.782} & 0.2500 & \textbf{4} \\
\bottomrule
\end{tabular}
\end{table}

The greedy router's zero-regret record is not evidence of good behaviour. It is an
artefact of a benchmark that degrades in lockstep with the damage being done: every
call greedy makes is optimal \emph{given the ecosystem it has already ruined}.

This paper formalises that failure. Our contributions:

\begin{itemize}
\item \textbf{Theorem~\ref{thm:vacuity}} (main result). A zero-regret policy can
incur value loss $m\delta E$, unbounded in pool size and externality magnitude, and
can finish net-negative while the optimum is positive. No-regret is not no-harm.
\item \textbf{Theorem~\ref{thm:char}}. Greedy is optimal if and only if
$\kappa_a = p_a e_a$ is constant across arms. Neither factor alone is the invariant.
\item \textbf{Theorem~\ref{thm:floor}}. The externality cannot be estimated to
arbitrary precision: the error floor is horizon-independent, because the
data-generating process is the consumption itself.
\item \textbf{Theorem~\ref{thm:seq}}. In the pure-sequencing regime the optimal
order is ascending in $e$; arm values are irrelevant.
\item A closed-form correction (ECI), a rollout policy reaching 99.5--99.9\% of
optimum, and an open-source testbed with one regression test per theorem.
\end{itemize}

\section{Related work}

\textbf{Rotting bandits} \cite{levine2017,seznec2019,seznec2020} model rewards that
decay with the number of pulls (\emph{rested}) or with elapsed time
(\emph{restless}); Seznec et al.~\cite{seznec2020} give a single adaptive-window
policy near-optimal for both. Decay there is gradual rather than removal, and
critically the learner is assumed to know which regime it is in.
\textbf{Mortal bandits} \cite{chakrabarti2008} have arms with stochastic lifetimes,
but expiry is exogenous. \textbf{Blocking bandits} make a pulled arm temporarily
unavailable. \textbf{Sleeping bandits} \cite{kleinberg2010} allow arbitrary
exogenous availability. None of these couple arms: destroying one does not change
another's reward distribution.

\textbf{Bandits with knapsacks} \cite{badanidiyuru2013} is the closest structural
neighbour: playing an arm consumes resources and the process halts when a budget is
exhausted. The budget is global and exogenous, and consumption leaves other arms'
rewards untouched. Here consumption permanently degrades \emph{every} future
reward. That difference is why the governing quantity is $\kappa = p\cdot e$ rather
than a consumption rate.

\textbf{Bandit learning with positive externalities}
\cite{shah2018} is the only prior work we are aware of that places externalities in
a bandit model. There, satisfied users attract similar users, so externalities are
positive and act on the \emph{arrival process}; the authors show UCB incurs linear
regret and develop balanced exploration. Our externalities are negative, act on
\emph{rewards}, and are triggered by irreversible destruction. The parallel is
worth drawing: in both settings a standard algorithm fails badly, indicating that
externality structure of either sign breaks classical guarantees.

Work on safe and conservative bandits constrains reward or maintains a baseline
while leaving arms available; work on open multi-agent systems models exogenous
population churn. Neither addresses action-triggered removal with reward coupling.

\section{Setting}

Arms $1,\dots,n$. Arm $a$ has mean reward $v_a$, destruction probability
$p_a \in (0,1]$, and externality coefficient $e_a \ge 0$. Let $S_t$ be the set of
available arms at round $t$ and define the accumulated burden
\[
B(S) \;=\; \delta \sum_{i \notin S} e_i .
\]
Pulling $a \in S_t$ yields reward $v_a - B(S_t) + \eta_t$ with $\eta_t$
sub-Gaussian with parameter $\sigma$, and destroys $a$ with probability $p_a$,
permanently adding $\delta e_a$ to the burden. The horizon is $T$. The value
function is
\[
V(S,t) = \max_{a \in S}\Big[ v_a - B(S) + p_a V(S\setminus\{a\}, t{+}1)
  + (1-p_a) V(S, t{+}1) \Big].
\]
Define the \emph{expected marginal externality}
\[
\boxed{\;\kappa_a \;=\; p_a\, e_a\;}
\]
the expected permanent burden created by one pull of $a$.

\textbf{Benchmark.} We measure regret against the best available arm,
$\max_{a \in S_t}\big(v_a - B(S_t)\big)$. This is the benchmark deployed systems
use, and Theorem~\ref{thm:vacuity} is a statement about that practice. Value is
measured against the exact optimum $V(S_0,0)$, computed by dynamic programming over
subsets for $n \le 8$ and by rollout above that.

\section{When is greedy optimal?}

\begin{theorem}[Characterisation]
\label{thm:char}
Greedy---pull $\arg\max_a v_a$---is optimal if and only if $\kappa_a$ is constant
across arms.
\end{theorem}

\begin{proof}[Proof sketch]
Since $B(S)$ and $V(S,t{+}1)$ do not depend on the chosen arm, maximising the
Bellman expression is equivalent to maximising $v_a - p_a D_a$ with
$D_a = V(S,t{+}1) - V(S\setminus\{a\},t{+}1)$. Decompose
$D_a = \delta e_a L(S,t) + O_a(S,t)$ into a burden and an option component. The
burden component contributes $p_a \delta e_a L = \delta \kappa_a L$, which is
identical across arms exactly when $\kappa$ is constant and hence drops out of the
argmax.

\emph{Sufficiency.} Under constant $\kappa$, couple the destruction randomness so a
pair $\{a,b\}$ receives the same two draws in either order. The surviving-set
distribution after both pulls is then identical, and since expected burden per pull
is $\delta\kappa$ for both arms, the expected burden paths coincide. The only
residual difference is which reward arrives first, giving $v_a - v_b \ge 0$.
Iterating sorts any optimal policy into greedy order.

\emph{Necessity.} The construction of Theorem~\ref{thm:vacuity} has
$\kappa_H = E > 0 = \kappa_S$ and a strictly positive gap $m\delta E$.
\end{proof}

Two empirical checks. Varying $p$ and $e$ widely while pinning $\kappa$ constant
leaves greedy exactly optimal (0/10 instances broken, gap $0.000\%$); varying the
product breaks it (10/10, gap $5.4$--$6.3\%$). Separately, the interchange
inequality underlying sufficiency was audited exhaustively over all reachable
states: $0$ violations in $6{,}842$ ordered pairs under constant $\kappa$, versus
$56.2\%$ violations when $\kappa$ varies. The inequality holds precisely when the
theorem says greedy is optimal.

The optimality gap grows monotonically with $\mathrm{std}(\kappa)$, from $5.9\%$ to
$29.1\%$ across six bins over 240 instances.

\section{A zero-regret policy can destroy unbounded value}

\begin{theorem}[Vacuity of no-regret under consumption]
\label{thm:vacuity}
For any $M > 0$ there is an instance on which greedy attains regret exactly $0$
against the best-available benchmark while $V^\star - V_{\mathrm{greedy}} \ge M$.
Moreover the loss can exceed $V^\star$, making the zero-regret policy net-negative
while the optimum is positive.
\end{theorem}

\begin{proof}
Fix $m \ge 1$, $E > 0$, $0 < \epsilon < \delta E$. Take $n = m{+}1$ arms: one
\emph{hot} arm $H$ with $v_H = 1$, $e_H = E$, and $m$ \emph{safe} arms with
$v_S = 1-\epsilon$, $e_S = 0$. Set $p_a = 1$ for all arms and $T = n$, so each arm
is pulled exactly once and only the order is chosen.

Greedy pulls $H$ first since $1 > 1-\epsilon$. $H$ is destroyed, adding permanent
burden $\delta E$, and each of the $m$ subsequent pulls yields
$(1-\epsilon) - \delta E$:
\[
V_{\mathrm{greedy}} = 1 + m(1 - \epsilon - \delta E).
\]
Deferring $H$ to last yields $1-\epsilon$ on each safe pull with no burden, and $1$
on the final pull of $H$, whose burden is charged only to later rounds, of which
there are none:
\[
V^\star \ge m(1-\epsilon) + 1 .
\]
Subtracting, $V^\star - V_{\mathrm{greedy}} = m\,\delta E$, unbounded in $m$ and in
$\delta E$. Choosing $\delta E > 1$ gives $V_{\mathrm{greedy}} < 0 < V^\star$.

At every round greedy pulls $\arg\max_{a\in S} (v_a - B(S))$, which is the
definition of the benchmark; its instantaneous regret is $0$ at every state.
\end{proof}

The closed form matches exact dynamic programming to machine precision at all $15$
settings tested, $m \in \{1,2,4,8,12\}$ and $\delta E \in \{0.5,1,2\}$. At $m=12$,
$\delta E = 2$: $V^\star = 12.88$, $V_{\mathrm{greedy}} = -11.12$, gap exactly
$24.0000$, regret $0.0000$.

\paragraph{A correction, recorded.} An earlier version of this construction claimed
the gap grows linearly in $T$. That is false: with $p=1$ the pool is exhausted
after $n$ pulls, so $T$ beyond $n$ is irrelevant. Under $p<1$ the gap grows with
$T$ only until $T \approx \sum_j 1/p_j$, then saturates---the same sample-budget
ceiling that appears independently in Theorem~\ref{thm:floor}. The correct scaling
is linear in pool size, which is the meaningful quantity: greedy's error is paid
once per subsequent pull.

\paragraph{Interpretation.} Greedy maximises $v_a$, the arm's \emph{private} value.
The optimal policy accounts for $\delta\kappa_a(T-t)$, the damage imposed on
everything else. Regret is a private-value metric; under externalities private and
social value diverge; therefore no-regret guarantees nothing about social outcome.

\section{The externality cannot be learned}

\begin{theorem}[Estimation floor]
\label{thm:floor}
For any estimator of $\delta e_i$,
\[
\mathrm{SE}(\delta \hat e_i) \;\ge\; \frac{2\sigma}{\sqrt{\,\min\!\big(T,\ \sum_j 1/p_j\big)}},
\]
which for $T > \sum_j 1/p_j$ is independent of the horizon.
\end{theorem}

\begin{proof}
All information about $e_i$ enters through the level shift $\delta e_i$ in the
reward process at the moment arm $i$ is destroyed. Conditioning on the destruction
time with $n_b$ rounds before and $n_a$ after, the MLE is the difference of pre-
and post-destruction sample means, with variance
$\sigma^2(1/n_b + 1/n_a)$. By AM--HM,
$1/n_b + 1/n_a \ge 4/(n_b+n_a) = 4/N$, with equality iff $n_b = n_a$, giving
$\mathrm{SE} \ge 2\sigma/\sqrt{N}$. Arm $j$ survives $\mathrm{Geometric}(p_j)$
pulls, so $\mathbb{E}[N] = \sum_j 1/p_j$, capped at $T$.
\end{proof}

All three components verify numerically: the variance formula matches Monte Carlo
to within $0.1\%$; the AM--HM floor is attained exactly at the balanced split
(ratio $1.0000$); and $\mathbb{E}[N] = \min(T, \sum_j 1/p_j)$ holds including when
$T$ binds.

The horizon-independence is stark. At $p = 0.9$, increasing $T$ from $20$ to $320$
changes RMSE by \emph{zero}---$0.2329$ at every horizon, because the pool caps the
data at $8.8$ rounds. At $p = 0.1$ the round count climbs $20 \to 40 \to 71 \to
82.3 \to 83.2$, flattening exactly at the predicted $80$.

\begin{corollary}[The central tension]
As $p$ rises, $\kappa = p e$ grows, so the externality matters more; simultaneously
$N \approx n/p$ shrinks and $\mathrm{SE} \ge 2\sigma\sqrt{p/n}$ degrades. The
parameter governing optimal behaviour becomes less estimable exactly as it becomes
more important.
\end{corollary}

This is not a sample-complexity result that more data resolves. The data-generating
process is the consumption itself: observing arm $i$'s externality requires
destroying arm $i$.

\paragraph{Structure does not rescue it.} Suppose $e_a = \langle \psi, z_a\rangle$
for known features $z_a \in \mathbb{R}^k$, $k \ll n$, so each destruction informs a
shared $k$-vector. Per-arm least squares is exactly worthless (identical to greedy
to three decimals at every coupling level); feature-based estimation recovers only
$17$--$28\%$ of the available gap; and a \emph{conservative} policy that assumes a
uniform upper bound and learns nothing outperforms both at every level. If the
mechanism were sample-sharing, capture should rise as $k$ falls. It does not:
capture is $4.5\%$ at $k=1$ and peaks at $30.3\%$ at $k=5$. At $k=1$ the
externality is nearly uniform, so $\kappa$ has little dispersion and by
Theorem~\ref{thm:char} there is nothing to capture. The obstacle is dispersion, not
sample size: the externality must vary for the correction to matter, and variation
is what makes it hard to estimate.

\section{Corrections}

\subsection{The externality-corrected index}

Theorem~\ref{thm:char} hands over the correction. Pulling $a$ creates expected
permanent burden $\delta\kappa_a$ charged against every remaining round:
\[
I(a,t) \;=\; v_a \;-\; \delta\,\kappa_a\,(T-t).
\]
Because $B(S)$ is common to all arms it drops out of the ranking, so ECI is a true
index: each arm's score depends only on its own parameters and global time, and the
ranking is invariant to accumulated damage.

Against exact DP, ECI captures $52\%$--$99\%$ of the optimality gap, with capture
improving as coupling strengthens. Its own gap is flat at $1.9$--$3.1\%$ across the
entire dispersion range while greedy's grows five-fold, so ECI's advantage grows
without bound in $\mathrm{std}(\kappa)$.

\subsection{Rollout}

No closed form exists for the option-loss term ECI omits. One step of policy
improvement over ECI captures it implicitly, since
$V_{\pi_0}(S) - V_{\pi_0}(S\setminus\{a\})$ is exactly the marginal value of
holding $a$. This reaches $99.5$--$99.9\%$ of exact optimum, and with Monte Carlo
policy evaluation runs at $n=40$ where exact DP is infeasible, adding $3.3\%$ over
ECI at strong coupling.

\subsection{An exactly solvable case}

\begin{theorem}[Pure sequencing]
\label{thm:seq}
If $p_a = 1$ for all $a$ and $T = n$, the optimal order is ascending in $e$, and
arm values are irrelevant to the optimal order.
\end{theorem}

\begin{proof}
Every arm is pulled exactly once, so $\sum_a v_a$ is a constant. Total value is
$\sum_a v_a - \delta \sum_s e_{a_s}(n-s)$: an arm at position $s$ pays its
externality $(n-s)$ times. Minimising requires pairing large $e$ with small $(n-s)$.
\end{proof}

Verified by brute force over all $n!$ orderings, exact at 8/8 instances. Notably ECI
is \emph{not} exactly optimal even here---a conjecture we falsified---and the exact
rule inverts the usual intuition: when consumption is certain, an arm's value is
irrelevant to scheduling; only its externality matters.

\section{Experiments}

Table~\ref{tab:main} compares policies on $20$ arms, $60$ rounds, $25$ seeds.

\begin{table}[t]
\centering
\caption{The regret and value columns are ordered backwards relative to each other.}
\label{tab:main}
\begin{tabular}{lrr}
\toprule
policy & value & regret \\
\midrule
greedy & 14.449 & \textbf{0.000} \\
Thompson sampling & 15.171 & 1.341 \\
UCB & 13.257 & 6.136 \\
conservative & 17.089 & 8.160 \\
ECI & \textbf{21.214} & 10.069 \\
\bottomrule
\end{tabular}
\end{table}

Thompson sampling is consistently \emph{worse} than greedy under known parameters.
Its exploration is not free here: sampling a suboptimal arm can destroy a
low-value, high-externality arm. At strong coupling the corrected index nearly
triples Thompson's realised value ($20.8$ vs.\ $7.2$). The mechanisms compose
rather than compete---TS handles uncertainty in $v$, ECI handles consumption---and
TS+ECI dominates both.

\section{Limitations}

The burden model $B(S) = \delta\sum e_i$ is additive and the simplest coupling that
exhibits the phenomenon; richer couplings (multiplicative, saturating, networked)
are unexplored. Exact optimality results rest on dynamic programming at $n \le 8$;
larger-scale claims rest on rollout. Theorem~\ref{thm:vacuity} is a statement about
the best-available-arm benchmark specifically---a forward-looking benchmark would
not be fooled, but it is also not what deployed systems measure, which is the point.
Theorem~\ref{thm:char}'s sufficiency proof relies on a coupling argument whose
formalisation under horizon truncation we verify exhaustively rather than prove in
full generality.

\section{Conclusion}

When actions consume the action set, regret is measured against an optimum the
learner has already degraded. A policy can hold regret at exactly zero while
driving value arbitrarily negative, and the parameter that would let it do better
becomes unlearnable precisely as it becomes important, because acquiring each
observation destroys the resource being priced. The practical consequence is
direct: for irreversible actions, externality costs must be \emph{specified}, not
discovered. A crude uniform upper bound outperforms every learned estimator we
tested.

\paragraph{Code.} Library, all 24 experiments including the falsified ones, and one
regression test per theorem: \url{https://github.com/<user>/consumable-bandits}

\begin{thebibliography}{9}
\bibitem{badanidiyuru2013} A.~Badanidiyuru, R.~Kleinberg, A.~Slivkins.
  Bandits with knapsacks. \emph{FOCS}, 2013.
\bibitem{chakrabarti2008} D.~Chakrabarti, R.~Kumar, F.~Radlinski, E.~Upfal.
  Mortal multi-armed bandits. \emph{NeurIPS}, 2008.
\bibitem{kleinberg2010} R.~Kleinberg, A.~Niculescu-Mizil, Y.~Sharma.
  Regret bounds for sleeping experts and bandits. \emph{Machine Learning}, 2010.
\bibitem{levine2017} N.~Levine, K.~Crammer, S.~Mannor.
  Rotting bandits. \emph{NeurIPS}, 2017.
\bibitem{seznec2019} J.~Seznec, A.~Locatelli, A.~Carpentier, A.~Lazaric,
  M.~Valko. Rotting bandits are no harder than stochastic ones.
  \emph{AISTATS}, 2019.
\bibitem{seznec2020} J.~Seznec, P.~Menard, A.~Lazaric, M.~Valko.
  A single algorithm for both restless and rested rotting bandits.
  \emph{AISTATS}, 2020.
\bibitem{shah2018} V.~Shah, R.~Johari, J.~Scheinkman.
  Bandit learning with positive externalities. \emph{NeurIPS}, 2018.
\end{thebibliography}

\end{document}
